\subsection{The Irreducible Sector and Spectral Hierarchy}
  \label{sec:baryonic-sector}

  Within the spectral admissibility framework established in Chapter~\ref{sec:spectral-admissibility},
  the $r=3$ configuration is the first sector in which algebraic
  irreducibility becomes structurally possible within the admissible
  subspace.
  If no non-trivial projector simultaneously preserves admissibility and
  diagonalizes $\mathcal{A}_3$, the resulting mass scale cannot be
  interpreted as a superposition of lower-sector contributions.

  \subsubsection*{From Dynamical Selection to Global Invariants}

    In the $r=2$ sector, mass amplification arises from a dynamical
    selection rule acting within a reducible two-mode structure.
    The Pisano ratio $\varphi$ enters as a stability condition on the
    torsional generator restricted to the first multiplet.

    For $r=3$, such a two-mode reduction need not exist.
    If the torsional action is algebraically irreducible in
    $H_{\mathrm{adm}}^{(3)}$, no analogue of the Pisano selection rule
    applies, because no non-trivial invariant subspace remains on which
    a two-mode frequency selection could operate.
    The relevant observable must then be a global spectral invariant.

    A natural candidate is the torsional action
    \(
    \mathcal{A}(3)
    \)
    defined in Eq.~\eqref{eq:Aw-fredholm},
    or an invariant ratio extracted canonically from the spectrum of
    \(
    \Delta^{(0)}_G+\Omega_3
    \)
    restricted to $H_{\mathrm{adm}}^{(3)}$.
    Such an invariant should be independent of overall amplitude
    rescalings of $\Omega_3$ and stable under increasing spectral
    truncation.

  \subsubsection*{Mass Hierarchy as Spectral Scaling}

    If the $r=3$ sector is irreducible,
    the sector mass scale is controlled by the spectral distortion
    induced by $\Omega_3$.
    In analogy with the reducible sectors one expects
    \begin{equation}
      \frac{m_{(3)}}{m_{(1)}}
      \;\sim\;
      \sqrt{
        \frac{\lambda_{\mathrm{bind}}^{(3)}}
        {\lambda_{1}}
      }.
    \end{equation}

    Unlike the reducible case,
    this ratio would not be determined by a two-mode dynamical fixed point,
    but by the global spectral structure of the irreducible torsional
    sector.

  \subsubsection*{Programmatic Outlook}

    The structural task is therefore twofold:

    \begin{itemize}
      \item[(i)] determine whether $\mathcal{A}_3'$ is trivial within
      $H_{\mathrm{adm}}^{(3)}$,
      \item[(ii)] identify a canonical spectral invariant of
      $\Delta^{(0)}_G+\Omega_3$ that is stable under admissible
      deformations.
    \end{itemize}

    If such an invariant exists and does not depend on parameter tuning,
    the irreducible-sector mass scale would emerge as a structural
    consequence of spectral irreducibility rather than as
    phenomenological input.
